\specialsection{Постановка задачи}

В данной работе рассматривается задача рекомендации следующей культуры для конкретного сельскохозяйственного поля на
основе его исторической последовательности использования и базовых признаков поля и региона.

Практически задача может быть представлена так:

\begin{center}\(
(\mathrm{history}_1, \mathrm{history}_2, \mathrm{history}_3, x) \rightarrow \mathrm{target}
\)\end{center}

где:

\begin{itemize}
\tightlist
\item
  \texttt{history\_1}, \texttt{history\_2}, \texttt{history\_3} --- культуры в предыдущие три года;
\item
  \texttt{x} --- дополнительные признаки поля и региона;
\item
  \texttt{target} --- культура следующего года.
\end{itemize}

Таким образом, с формальной точки зрения задача сводится к многоклассовой классификации следующего шага в
последовательности культур. Однако прикладной смысл работы не ограничивается только top-1 предсказанием. В реальной
задаче выбора следующей культуры специалисту полезен не один ответ, а несколько правдоподобных вариантов, которые можно
дополнительно оценить с учётом внешних ограничений и предметного контекста. Поэтому в работе используется
рекомендательная интерпретация модельного результата: система формирует ранжированный список наиболее вероятных
кандидатов.

Входными данными для модели являются:

\begin{itemize}
\tightlist
\item
  история культур за три предыдущих года;
\item
  региональные признаки;
\item
  признаки, описывающие само поле.
\end{itemize}

В качестве таких признаков в работе используются:

\begin{itemize}
\tightlist
\item
  коды штата и округа;
\item
  площадь поля;
\item
  пространственные координаты.
\end{itemize}

Выход системы интерпретируется в двух уровнях:

\begin{itemize}
\tightlist
\item
  как наиболее вероятная следующая культура;
\item
  как Top-k shortlist вероятных кандидатов.
\end{itemize}

Логика работы в прикладном виде может быть представлена так:

\begin{center}\(
\text{история поля} + \text{признаки поля} \rightarrow \text{ранжированный список следующих культур}
\)\end{center}

В этой постановке решаются три взаимосвязанные подзадачи:

\begin{itemize}
\tightlist
\item
  \textbf{предсказание} следующей культуры по истории поля и дополнительным признакам;
\item
  \textbf{ранжирование} кандидатов и формирование короткого списка вероятных вариантов;
\item
  \textbf{интерпретация} полученного shortlist с учётом уверенности модели и объяснительного контекста.
\end{itemize}

Важно подчеркнуть, что работа не ставит целью построение полного оптимизатора севооборота. Разрабатываемая система
рассматривается как объяснимый рекомендательный прототип, который помогает анализировать вероятные варианты следующей
культуры, но не заменяет экспертное агрономическое решение.

В более компактной формулировке задача работы заключается в следующем: по истории культур на поле и базовым признакам
поля и региона требуется построить систему, которая формирует ранжированный список вероятных следующих культур,
оценивает уверенность рекомендаций и предоставляет интерпретируемый контекст для анализа shortlist-кандидатов.

Для решения этой задачи необходимо:

\begin{itemize}
\tightlist
\item
  подготовить и структурировать исторические данные по культурам;
\item
  сформировать обучающую выборку в виде временных окон;
\item
  построить предиктивное ядро модели;
\item
  оценить качество результата в Top-k постановке;
\item
  дополнить модель confidence- и explainability-слоями.
\end{itemize}

Такая постановка позволяет рассматривать выбор следующей культуры не только как задачу классификации, но и как задачу
поддержки принятия решений, в которой важны не только точность предсказания, но и удобство интерпретации результата.
